% TravelView - pandoc/xelatex 中文 PDF 支持
% 由 tools/pandoc/defaults.yaml 通过 include-in-header 引入
%
% 背景: pandoc 默认用 pdflatex + Helvetica，没有任何中文字形，
%       于是每个汉字都报 "Missing character"，PDF 里整片空白。
%       修复的三个要素: (1) 换 xelatex (2) 指定中文字体 (3) 让符号也走中文字体。

\usepackage{xeCJK}
\usepackage{newunicodechar}

% --- 中文字体 (macOS 自带 PingFang SC，无需安装) ---
\setCJKmainfont{PingFang SC}[BoldFont = PingFang SC Semibold]
\setCJKsansfont{PingFang SC}

% --- 代码块字体 (关键) ---
% 想让 ASCII 流程图框线严格对齐，须用「中文=西文两倍宽」的等宽字体
% (Sarasa Mono SC)。本机当前未安装，故先用 PingFang SC (自带)，代价是
% 代码里夹杂中文时流程图边界会略微错位。要完美对齐:
%     brew install --cask font-sarasa-gothic
%     再把下面一行换回  \setCJKmonofont{Sarasa Mono SC}
\setCJKmonofont{PingFang SC}

% --- 让制表符 / 箭头 / 带圈数字 也交给中文字体 ---
% Helvetica、Menlo 里没有这些字形，不做这一步同样会掉字。
\xeCJKDeclareCharClass{CJK}{
  "2010 -> "2027,   % 破折号、省略号
  "2190 -> "21FF,   % 箭头
  "2460 -> "24FF,   % 带圈数字
  "2500 -> "257F,   % 制表符 (ASCII 流程图边框)
  "25A0 -> "25FF,   % 几何图形
  "2600 -> "27BF    % 杂项符号
}

% emoji 的变体选择符 U+FE0F 直接吞掉，防止触发缺字警告
\newunicodechar{^^^^fe0f}{}

% 代码块自动换行，避免长命令溢出页面
\usepackage{fvextra}
\DefineVerbatimEnvironment{Highlighting}{Verbatim}{breaklines,breakanywhere,commandchars=\\\{\}}
